\documentclass[12pt]{article}
\usepackage[utf8]{inputenc}
\usepackage[T1]{fontenc}
\usepackage{hyperref}
\title{Tolomako Language Resources for K-12 Curriculum Development}
\author{Compiled by Claude Sonnet 5 (claude-sonnet-5) for Andr\'as Kornai}
\date{}
\begin{document}
\maketitle

\section{Introduction}
Tolomako (also Bigbay, ISO 639-3 tlm) is an Oceanic language of the Big Bay area of Espiritu Santo, Vanuatu. It has about 900 speakers by Wikipedia's count, and the UNESCO Atlas classes it as Definitely Endangered. Joshua Project gives higher population figures for the people (1,500 to 1,600). The language has no tense or aspect marking and distinguishes singular, dual, trial and plural pronouns. No school materials were found.

\section{Primary Educational Resources}
None found in this search.

\section{Linguistic and Reference Materials}
\begin{itemize}
\item Wikipedia article on the Tolomako language: \url{https://en.wikipedia.org/wiki/Tolomako_language}
\item Tolomako entry in the TLOPO database (Tolomako, Sakao and related languages): \url{https://tlopo.clld.org/languages/tolomako}
\item Glottolog map for Tolomako: \url{https://glottolog.org/resource/languoid/id/tolo1255.bigmap.html}
\item Rosetta Project page for Tolomako: \url{https://rosettapanglossia.longnow.org/index.php?title=Tlm}
\end{itemize}

\section{Religious and Community Resources}
\begin{itemize}
\item Joshua Project profile of the Tolomako people: \url{https://joshuaproject.net/people_groups/15543}
\end{itemize}

\section{Digital and Audio Resources}
None found beyond the database entries above.

\section{Curriculum Development Context}
Only descriptive and database materials turned up. No orthography, primers or scripture were seen. A curriculum would start from community-led orthography work and documentation. Because young people do not all speak it, a language-revitalization programme would be the first step.

\end{document}